\documentclass[a4paper, 12pt]{article}
% math symbols
\usepackage{amssymb}
\usepackage{amsmath}
\usepackage{mathrsfs}
\usepackage{mathseries}


\usepackage[margin = 2cm]{geometry}

\tolerance = 1000
\emergencystretch = 0.74cm



\pagestyle{empty}
\parindent = 0mm

\renewcommand{\coursetitle}{DM/ML}
\setcounter{curtask}{1}

\setmathstyle{\includegraphics[scale = 0.05]{../pics/utia-rest.png}}{Мат. логика}{3 курс}


\begin{document}

\task*{
    Три узника сидят каждый в своей камере. Надзиратель сообщает узникам, что на следующий день запустит
    один из двух процессов:
    \begin{itemize}
        \item для каждого натурального $n$ в момент времени $\frac{1}{n}$ секунды до полуночи равно в
            одной камере включается свет и сразу гаснет;
        \item для конечного количества натуральных $n$ происходит тоже, что и в первом случае, а для
            всех остальных моментов времени вида $\frac{1}{n}$ секунда до полуночи свет включается ровно
            в двух камерах.
    \end{itemize}
    В полночь каждый из узников объявляет какой процесс запускал наблюдатель. Если хотя бы двое угадали,
    то их всех отпускают. Можно ли гарантированно этого добиться? Узники могут заранее договориться и
    могут обрабатывать бесконечное число информации.
}

\begin{definition*}
    Будем называть тип $p$ теории $T$ \deftext{изолированным}, если найдется такое конечное множество
    формул $\varphi_1(v), \varphi_2(v), \dots, \varphi_k(v) \in p$, что для любой формулы $\psi(v) \in p$:
    $$
        T \vdash \forall v, \left(\bigwedge\limits_{i \le k} \varphi_i \rightarrow \psi(v)\right).
    $$
\end{definition*}

\task{
    Покажите, что любой изолированный тип полной теории $T$ реализуется в любой модели теории $T$.
}

\task{
    Предъявите пример полной теории $T$, у которой есть неизолированный полный $1$-тип.
}

\task{
    Покажите, что у теории плотных линейных порядков без первого и последнего элемента конечное число
    полных $n$-типов.
}

\task{
    Рассмотрим теорию $\Theory(\mathbb{Z}, s, <)$, $s(x) \coloneqq x + 1$. Какие из полных $n$-типов
    изолированы?
}

\task{
    Пусть $n \in \mathbb{N}$. Покажите, что:
    $\PrSys{MA} \vdash (n < x) \leftrightarrow ((0 \neq x) \land (s(0) \neq x) \land (s(s(0)) \neq x)
    \land \cdots \land (s^{(n)}(0) \neq x)$.
}

\task{
    Покажите, что:
    $\PrSys{MA} \vdash s^{(n)}(0) + s^{(m)}(0) = s^{(n + m)}(0)$.
}

\task{
    Покажите, что в $\PrSys{MA}$ нельзя вывести:
    $(x + y) + z = x + (y + z)$.
}





\end{document}



%%% Local Variables:
%%% mode: latex
%%% TeX-master: t
%%% End: